(1) 透热且无摩擦可动活塞在长期热与机械平衡下要求 \(T_{1f}=T_{2f}=T_f\) 及 \(p_{1f}=p_{2f}=p_f\)。由于外壳既不与外界换热也不对外做功，整个气体总内能保持不变；活塞移动造成的两室内功彼此抵消。
(2) 两室内能为 \(U_i=n_iC_VT_i\)，故末温 \(T_f=(n_1T_1+n_2T_2)/(n_1+n_2)\)。由 \(p_fV_{if}=n_iRT_f\) 与 \(V_{1f}+V_{2f}=V\)，得到 \(V_{1f}=n_1V/(n_1+n_2)\)、\(V_{2f}=n_2V/(n_1+n_2)\)；各气室的体积份额只由摩尔数决定。
(3) 熵增为 \(\Delta S=\sum_{i=1}^2 n_i[C_V\ln(T_f/T_i)+R\ln(V_{if}/V_{ii})]\)。以每室内能和容积为变量，总熵函数对正温、正容积严格凹，约束 \(U_1+U_2=U_{\rm tot},V_1+V_2=V\) 下唯一极大点由 \(\partial S/\partial U_i=1/T_i\) 与 \(\partial S/\partial V_i=p_i/T_i\) 的拉格朗日乘子条件给出共同 \(T,p\)；因此平衡熵不小于任意初态熵。
(4) 末压 \(p_f=(n_1+n_2)RT_f/V\)，末内能 \(U_f=(n_1+n_2)C_VT_f=n_1C_VT_1+n_2C_VT_2\)，核对守恒。若 \(n_1=n_2\)，两室末体积各为 \(V/2\)；若初态已经满足共同 \(T,p\)，则初体积必按摩尔数分配，所有末态等于初态，熵增严格为零。
